\chapter{Conclusion}\label{chapter:conclusion}

This survey reviewed KV-cache infrastructure for LLM inference through three responsibilities: representation and compression, storage and access, and reuse and execution management. We analyzed these techniques against capacity pressure, access bandwidth pressure, and state preparation overhead. This framework connects model-level changes with the systems that retain, move, and reuse their cached state, within and across servers.

The reviewed work shows that relieving one bottleneck can increase pressure on another. Compression reduces data volume but may introduce processing costs or quality loss. Offloading expands capacity while making access and restoration more expensive. Reuse avoids repeated computation, but retained state consumes storage and must be compatible with the execution that uses it. The benefit of a technique therefore depends on the model, workload, hardware, and other mechanisms used alongside it. Effective infrastructure requires coordinated choices across these responsibilities.

Future systems must support both ultra-long contexts and long-lived agentic tasks. Longer contexts require affordable access to extensive histories while preserving information needed later. Longer task lifetimes require inexpensive state preservation during tool waits and timely recovery when execution resumes. Together, these trends extend cache management from individual inference requests to complete tasks. Balancing capacity, access, and preparation costs while maintaining task quality will remain central to this development.
